\section{Auxiliary proofs}
\label{sec:auxiliary-proofs}

% Alternative treedepth proof retained in source but omitted from the paper.
% \subsection{The earlier incidence construction}
% \label{app:incidence-proof}
%
% We explain how the incidence construction of Wu, Austrin, Pitassi, and
% Liu~\cite[Sections~4.3 and~5]{APW} already implies the qualitative
% dense treedepth gap in Section~\ref{subsec:sse-dense-gap}. The inherited
% part is the source reduction and its separator soundness analysis.
% The additional argument is a treedepth upper bound in the YES case.
%
% \begin{restatedtheorem}{Incidence-construction consequence}
% For every fixed $C\geq1$, approximating treedepth within factor $C$
% is SSE-hard even on simple bipartite graphs.
% \end{restatedtheorem}
%
% \emph{The output graph.}
% Let $G$ be their regular source host, with $m>0$ edges, and let
% $R_{\mathrm{inc}}\geq |V(G)|m$ be an integer. Denote its incidence output by
% $G'$. Its vertex set is the disjoint union of $V(G)\times[R_{\mathrm{inc}}]$ and
% $E(G)$: $(v,s)$ is the $s$th copy of the host vertex $v$, whereas
% the output vertex indexed by $e\in E(G)$ represents a host edge.
% The adjacency rule is
% \[
%  \{e,(v,s)\}\in E(G')
%  \quad\Longleftrightarrow\quad
%  e\in E(G),\ v\in e,\ s\in[R_{\mathrm{inc}}].
% \]
% There are no edges within either $V(G)\times[R_{\mathrm{inc}}]$ or $E(G)$, and
% no other edges. In particular, $G'$ is simple and bipartite.
%
% \emph{The YES forest.}
% The source YES partition $V_1,\ldots,V_q$ has equal class volumes
% $\vol_G(V_i)=2m/q$ and $\Phi_G(V_i)\leq2\varepsilon$.
% Let $D$ be the multiset of crossing host edges and
% $E_i=E(G[V_i])$ the multiset of internal host edges in class $i$.
% Each crossing edge belongs to the cuts of both endpoint classes. Therefore
% \[
%  2|D|=\sum_{i=1}^q|\cut_G(V_i)|
%  =\sum_{i=1}^q\Phi_G(V_i)\vol_G(V_i)
%  \leq4\varepsilon m.
% \]
% Each internal edge contributes two incidences to its class volume,
% so $2|E_i|\leq\vol_G(V_i)=2m/q$. Consequently
% $|D|\leq2\varepsilon m$ and $|E_i|\leq m/q$ for every $i$.
% These are also the unweighted bounds of
% Lemma~\ref{lem:sse-yes-edge-counts}.
%
% First make a top chain from the output vertices indexed by $D$.
% For each $i$, make a chain from the output vertices indexed by
% $E_i$ and place it below the last vertex of the top chain. Attach
% all vertices $(v,s)$ with $v\in V_i$ as leaves below the last
% vertex of this class chain. Empty chains are omitted: if $E_i$ is
% empty, these leaves go directly below the top chain. If $D$ is
% empty, the class chains and any unattached leaves are separate
% rooted trees. This describes a rooted forest on every vertex of
% $G'$, including copies of isolated host vertices.
%
% To verify that it is an elimination forest, consider an arbitrary
% output edge $\{e,(v,s)\}$. If $e\in D$, its output vertex lies on
% the top chain and is an ancestor of $(v,s)$. Otherwise $e\in E_i$
% for some $i$, both host endpoints of $e$ lie in $V_i$, and the
% output vertex indexed by $e$ is on the class chain above $(v,s)$.
% Thus the endpoints of every output edge are comparable. A
% root-to-leaf path uses at most $|D|$ vertices of the top chain,
% at most $\max_i|E_i|$ vertices of one class chain, and one host
% vertex copy. Consequently
% \[
%  \td(G')\leq |D|+\max_i|E_i|+1
%  \leq 2\varepsilon m+\frac{m}{q}+1.
% \]
% The same bound covers empty top or class chains.
%
% \emph{The NO bound and the gap.}
% On a NO source instance, every set $A\subseteq V(G)$ with
% $|A|\in[|V(G)|/10,9|V(G)|/10]$ has at least
% $c_0\sqrt{\varepsilon}\,m$ crossing host edges, for an absolute
% constant $c_0>0$. If $d$ is the positive degree of the regular host,
% then $2m=d|V(G)|$ and
% $\mu_G(A)=d|A|/(d|V(G)|)=|A|/|V(G)|$. Hence this cardinality range
% is the same as the corresponding volume range. We use the following
% separator soundness of the incidence construction.
%
% \begin{lemma}[Incidence separator soundness
% {\cite[Lemma~5.1]{APW}}]
% \label{lem:incidence-separator-soundness}
% There is an absolute constant $c_1>0$ such that the replicated
% incidence graph $G'$ constructed above from a NO instance of the
% regular APW source satisfies
% \[
%  b(G')\geq c_1\sqrt{\varepsilon}\,m.
% \]
% The separator is balanced with respect to all vertices of $G'$,
% including the replicated host vertices.
% \end{lemma}
%
% Set $c_{\mathrm{inc}}=\min\{c_0,c_1\}$, with $c_1$ from
% Lemma~\ref{lem:incidence-separator-soundness}.
% Lemma~\ref{lem:balanced-by-td} therefore gives
% $\td(G')\geq c_{\mathrm{inc}}\sqrt{\varepsilon}\,m$.
%
% For completeness, fix any desired approximation factor $C\geq1$.
% Choose a sufficiently small fixed $0<\varepsilon\leq1/2$ for which
% the source promise applies and $\sqrt{\varepsilon}<c_{\mathrm{inc}}/(4C)$, and
% put $q=\lceil1/\varepsilon^2\rceil$. In the YES case, all $q$
% partition classes have positive volume and are therefore nonempty.
% The regular host has positive integer degree, so
% $2m\geq|V(G)|\geq q\geq1/\varepsilon^2$.
% Since $\varepsilon\leq1/2$, this implies
% $m\geq1/(2\varepsilon^2)\geq1/\varepsilon$. Also,
% $q\geq\varepsilon^{-2}$ and $m\geq\varepsilon^{-1}$ give
% $m/q\leq\varepsilon^2m$ and $1\leq\varepsilon m$. Thus the YES
% upper bound satisfies
% \[
%  2\varepsilon m+\frac{m}{q}+1
%  \leq(2\varepsilon+\varepsilon^2+\varepsilon)m
%  \leq4\varepsilon m,
% \]
% where the last inequality uses $\varepsilon^2\leq\varepsilon$.
% The NO lower bound is
% $c_{\mathrm{inc}}\sqrt{\varepsilon}\,m>4C\varepsilon m$. A factor-$C$
% approximation returns height at most $4C\varepsilon m$ in the YES
% case, while every feasible forest has greater height in the NO case.
% Thus it distinguishes the two cases. Finally, taking
% $R_{\mathrm{inc}}=|V(G)|m$ makes the incidence output polynomial in the source
% size. This proves SSE hardness for treedepth at every fixed factor
% on bipartite outputs from the cited construction.

\subsection{The branching inequality}
\label{app:branching-proof}

\begin{restatedlemma}{Lemma~\ref{lem:branching-ineq}}
For a positive integer $\nu$, positive reals $q_1,\ldots,q_\nu$, and
nonnegative reals $n_1,\ldots,n_\nu$,
\[
  3\sum_{j=1}^\nu q_jn_j-2\max_{1\leq j\leq\nu}(q_jn_j)
  \leq\left(\sum_{j=1}^\nu q_j\right)
        \left(\sum_{j=1}^\nu n_j\right).
\]
\end{restatedlemma}

\begin{proof}[Proof of Lemma~\ref{lem:branching-ineq}]
Put $x_j=\sqrt{q_jn_j}$ for $1\leq j\leq\nu$, and choose
$k\in\{1,\ldots,\nu\}$ with $x_k=\max_{1\leq j\leq\nu}x_j$.
Cauchy--Schwarz gives the first inequality below. All cross products
are nonnegative, and $x_kx_j\geq x_j^2$ because
$0\leq x_j\leq x_k$. Therefore
\[
\begin{gathered}
 \left(\sum_{j=1}^{\nu}q_j\right)
 \left(\sum_{j=1}^{\nu}n_j\right)
 \overset{\text{Cauchy--Schwarz}}{\ge}
   \left(\sum_{j=1}^{\nu}x_j\right)^2 =\sum_{j=1}^{\nu}x_j^2
   +2\sum_{1\leq i<j\leq\nu}x_ix_j
 \ge \sum_{j=1}^{\nu}x_j^2
       +2x_k\sum_{1\leq j\leq\nu,\;j\ne k}x_j\\
 \ge \sum_{j=1}^{\nu}x_j^2
       +2\sum_{1\leq j\leq\nu,\;j\ne k}x_j^2
 =3\sum_{j=1}^{\nu}x_j^2-2x_k^2.
\end{gathered}
\]
Substituting $x_j^2=q_jn_j$ proves the claim.
\end{proof}

\subsection{The quadratic branching inequality}
\label{app:total-depth-branching}

\begin{restatedlemma}{Lemma~\ref{lem:quadratic-branching}}
For a positive integer $\nu$ and nonnegative reals
$W_1,\ldots,W_\nu$,
\[
 \sum_{j=1}^\nu W_j^2+14\sum_{1\leq i<j\leq\nu}W_iW_j
 \leq\left(\sum_{j=1}^\nu\sqrt{W_j}\right)^4.
\]
\end{restatedlemma}

\begin{proof}[Proof of Lemma~\ref{lem:quadratic-branching}]
Put $x_j=\sqrt{W_j}\geq0$ for $1\leq j\leq\nu$. The fourth-power
expansion has pure terms $x_j^4$ and, for each pair $1\leq i<j\leq\nu$, the three mixed terms
$4x_i^3x_j+6x_i^2x_j^2+4x_ix_j^3$. All terms involving at least
three distinct indices are nonnegative. For each pair,
\[
4x_i^3x_j+6x_i^2x_j^2+4x_ix_j^3
 =14x_i^2x_j^2+4x_ix_j(x_i-x_j)^2
 \geq14x_i^2x_j^2.
\]
Sum these bounds, add the pure terms, and discard the terms with
three or more indices. Substituting $x_j^2=W_j$ proves the claim.
\end{proof}

\subsection{Derivation of the anchor-length corollary}
\label{app:anchor-proof}

We prove Corollary~\ref{cor:max-part-anchor} for the anchor graph
$H_\ell(G)$, its cliques $K_X$, and vertices $z_X$ from
Section~\ref{subsubsec:anchor-graph}. Fix a tripartite graph $G$ with
nonempty specified independent parts $A,B,C$ and an integer $\ell\ge1$.
Put
\[
 \vc_2(G)=\min_{X,Y\in\{A,B,C\},\;X\ne Y}\vc(G[X\cup Y]).
\]
The following published identity supplies the case of long anchors.

\begin{lemma}[BNS anchor identity {\cite[Lemma~6]{BNS}}]
\label{lem:bns-anchor}
Let $G$ be a tripartite graph with specified nonempty independent parts
$A,B,C$. For every positive integer $L\ge\vc(G)$,
$\td(H_L(G))=L+\vc(G)+1$.
\end{lemma}

For an elimination forest $F$ of $H_\ell(G)$, write $h(F)$ for its
height and define, for $X\in\{A,B,C\}$,
\[
 v_X=\underset{v\in K_X}{\operatorname{arg\,max}}\,
                   \operatorname{depth}_F(v).
\]
This vertex is unique because the vertices of a clique form an
ancestor chain.

\begin{lemma}[A path through an anchor {\cite[Claim~7 in the proof of Lemma~6]{BNS}}]
\label{lem:anchor-path}
Let $X\in\{A,B,C\}$. If every vertex of
$R\subseteq V(H_\ell(G))\setminus(K_X\cup\{z_X\})$ is ancestral
to $v_X$, then
\begin{equation}\label{eq:anchor-path-bound}
 h(F)\ge |R|+\ell+1.
\end{equation}
\end{lemma}

\begin{proof}
Since $K_X$ is a clique, all its vertices are ancestral to its deepest
vertex $v_X$. Thus $R\cup K_X$ lies on the ancestor path ending at
$v_X$. The edge $z_Xv_X$ makes $z_X$ comparable with $v_X$, so
$R\cup K_X\cup\{z_X\}$ lies on one root-to-leaf path. The sets
are disjoint and contain $|R|+\ell+1$ vertices in total, proving
the bound. This argument applies to every elimination forest $F$
and every $\ell\ge1$.
\end{proof}

The next lemma adapts the argument of
\cite[Claim~9 in the proof of Lemma~6]{BNS}. Only the two specified anchors need to be
incomparable; no condition on the third anchor is required.

\begin{lemma}[A cover from two incomparable anchors]
\label{lem:anchor-pair-cover}
Let $X,Y\in\{A,B,C\}$ be distinct and define
\[
 S_{X,Y}=\{v\in X\cup Y:v\text{ is ancestral to both }v_X
                                      \text{ and }v_Y\text{ in }F\}.
\]
If $v_X,v_Y$ are incomparable, then $S_{X,Y}$ is a vertex cover
of $G[X\cup Y]$.
\end{lemma}

\begin{proof}
For an edge $xy\in E(G[X\cup Y])$ with $x\in X$ and $y\in Y$,
the four vertices $v_X,x,y,v_Y$ induce a connected subgraph, since
$v_Xx$, $xy$, and $yv_Y$ are edges of $H_\ell(G)$.
Lemma~\ref{lem:connected-set-ancestor} gives one of these four
vertices ancestral to all four. It cannot
be $v_X$ or $v_Y$, since those anchors are incomparable. Hence $x$
or $y$ is ancestral to both anchors. Thus
$\{x,y\}\cap S_{X,Y}\ne\varnothing$ for every
$xy\in E(G[X\cup Y])$, as required.
\end{proof}

\begin{lemma}[Extending incomparable anchors]\label{lem:incomparable-anchor-extension}
If $v_A,v_B,v_C$ are pairwise incomparable, then
$h(F)\ge\ell+\vc(G)+1$.
\end{lemma}

\begin{proof}
Set $L=\max\{\ell,\vc(G)\}$ and
$d=L-\ell$. In each $K_X$, insert a chain of $d$ new vertices
immediately above $v_X$. Together with $K_X$, these vertices form
the size-$L$ anchor clique in $H_L(G)$, which is joined completely
to $X\cup\{z_X\}$. Every vertex of $X\cup\{z_X\}$ is comparable
with $v_X$, since it is adjacent to $v_X$. Hence it is
also comparable with the inserted chain. The resulting forest is
therefore an elimination forest of $H_L(G)$. Since the three
anchors are pairwise incomparable, each root-to-leaf path meets
at most one inserted chain, so its height is at most $h(F)+d$.
We now have
$h(F)+d\overset{\text{Lem.~\ref{lem:bns-anchor}}}{\ge}L+\vc(G)+1$, and thus
$h(F)\ge L+\vc(G)+1-d=\ell+\vc(G)+1$.
\end{proof}

We now distinguish the possible ancestor relations among the three
anchors to obtain a bound for every anchor length.

\begin{lemma}[Bounds for arbitrary anchor length]
\label{lem:arbitrary-anchor}
For every tripartite graph $G$ with nonempty specified independent
parts $A,B,C$ and every integer $\ell\ge1$,
\[
 1+\min\{\ell+\vc(G),\ 2\ell+\vc_2(G),\ 3\ell\}
 \le\td(H_\ell(G))\le\ell+\vc(G)+1.
\]
\end{lemma}

\begin{proof}
\emph{Upper bound.}
We follow the cover-first forest in the proof of~\cite[Lemma~6]{BNS}.
Choose a vertex cover $S\subseteq V(G)$ with $|S|=\vc(G)$,
place $S$ on a top chain, then place each anchor clique
on its own chain below it and attach the uncovered vertices of
that part and its associated $z_X$ as leaves below the anchor chain.
If the cover is empty, use the three anchor chains as separate roots.
Every source edge meets the top cover, and every anchor edge has
comparable endpoints. A root-to-leaf path therefore uses at most
$\vc(G)+\ell+1$ vertices. The construction does not require
$\ell\ge\vc(G)$, so it gives the upper bound for every $\ell\ge1$.

\emph{Lower bound.}
Fix an elimination forest $F$ and use the anchor notation above.
There are three cases.
\begin{enumerate}
 \item If $v_A,v_B,v_C$ are pairwise incomparable,
   $h(F)\overset{\text{Lem.~\ref{lem:incomparable-anchor-extension}}}{\ge}
   \ell+\vc(G)+1$.

  \item Suppose all three anchors are comparable, and let $v_X$
    be the deepest, with $\{Y,Z\}=\{A,B,C\}\setminus\{X\}$.
    Every vertex of the other two anchor cliques lies above $v_X$.
    Set $R=K_Y\cup K_Z$, so $|R|=2\ell$.
    Then
    $h(F)\overset{\eqref{eq:anchor-path-bound}}{\geq}|R|+\ell+1=3\ell+1$.

 \item Otherwise some pair is comparable, but the three anchors do
   not form a chain. Relabel so that $v_A$ is an ancestor of
   $v_B$ and $v_B$ is incomparable with $v_C$.
   This includes the case where $v_A$ is above both other
   anchors. Set $S=S_{B,C}$. The cover bound is
   \begin{equation}\label{eq:anchor-common-cover-bound}
    |S|\overset{\text{Lem.~\ref{lem:anchor-pair-cover}}}{\ge}
    \vc(G[B\cup C])\ge\vc_2(G).
   \end{equation}
   Every vertex
   of $K_A$ is an ancestor of $v_A$, and hence of $v_B$.
   Therefore $R=S\cup K_A$ is disjoint from $K_B\cup\{z_B\}$ and
    consists of ancestors of $v_B$. Since $S\subseteq V(G)$ and
    $K_A$ is disjoint from $V(G)$, $|R|=|S|+|K_A|=|S|+\ell$.
    Applying the path bound with anchor $B$ gives
    \[
    h(F)\overset{\eqref{eq:anchor-path-bound}}{\geq}|R|+\ell+1
           =2\ell+|S|+1
  \overset{\eqref{eq:anchor-common-cover-bound}}{\geq}
           2\ell+\vc_2(G)+1.
    \]
\end{enumerate}
Taking the weakest of these three bounds proves the lemma.
\end{proof}

\begin{restatedcorollary}{Corollary~\ref{cor:max-part-anchor}}
If $G$ is tripartite with specified nonempty independent parts $A,B,C$
and $\ell\geq\max\{|A|,|B|,|C|\}$, then
$\td(H_\ell(G))=\ell+\vc(G)+1$.
\end{restatedcorollary}

\begin{proof}[Proof of Corollary~\ref{cor:max-part-anchor}]
Since each part has size at most $\ell$, any two parts together
cover every edge and $\vc(G)\le2\ell$. Choose a pair $X,Y$
attaining $\vc_2(G)$ and let $Z$ be the omitted part. A minimum
vertex cover of $G[X\cup Y]$, together with all of $Z$, covers $G$.
Therefore $\vc(G)\le\vc_2(G)+|Z|\le\vc_2(G)+\ell$. The two
cover bounds give
$\ell+\vc(G)\le3\ell$ and
$\ell+\vc(G)\le2\ell+\vc_2(G)$, respectively. Thus the minimum
is $\ell+\vc(G)$, and
\[
 \td(H_\ell(G))
 \overset{\text{Lem.~\ref{lem:arbitrary-anchor}}}{=}
 \ell+\vc(G)+1.
\]
\end{proof}

\subsection{Total-depth normalization and the three paths}
\label{app:total-depth-tools}

Let $G$ be a simple tripartite graph with specified nonempty
independent parts $A,B,C$, and put $N=|V(G)|$. Complete each part
to a clique by setting
\[
 \widehat G=\left(V(G),\ E(G)\cup
                 \bigcup_{X\in\{A,B,C\}}\binom{X}{2}\right).
\]
For an integer $M\ge1$, obtain
$H^\Sigma_M(G)$ from $\widehat G$ by adding disjoint independent
sets $T_A,T_B,T_C$, with
\[
 |T_X|=M,\qquad N_{H^\Sigma_M(G)}(t)=X
 \qquad(X\in\{A,B,C\},\ t\in T_X).
\]
Depths in a rooted
forest start at one, and ancestor includes the vertex itself. An
elimination forest has one node per graph vertex and makes the
endpoints of every edge comparable in the ancestor order.
Write $\totd(F)$ for the sum of vertex depths in $F$, and
$\totd(J)$ for its minimum over elimination forests of a graph $J$.

In any elimination forest $F_G$ of $\widehat G$, the vertices of a
part are pairwise comparable because they form a clique. They lie
on one ancestor path and have a unique deepest vertex $y_X$. Define
\[
 W(F_G)=\sum_{X\in\{A,B,C\}}
             \bigl(\operatorname{depth}_{F_G}(y_X)+1\bigr),
 \qquad W(G)=\min_{F_G}W(F_G).
\]
Here and below, minima over $F_G$ range over all elimination forests
of $\widehat G$.
For a choice of part $X$, let $Y,Z$ be the other two parts. Write
$\vc(J)$ for the minimum cardinality of a vertex cover of $J$, and
define
\[
 \operatorname{cost}_X(G)=\min_{S_0}
 \left(2|S_0|+\vc\bigl(G[(Y\cup Z)\setminus S_0]\bigr)\right),
\]
where $S_0\subseteq V(G)$ covers every $X$--$(Y\cup Z)$ edge.
All cover graphs below refer to $G$, before clique completion.

\begin{lemma}[Moving one added vertex]
\label{lem:added-vertex-move}
For any $X\in\{A,B,C\}$ and $u\in T_X$, every elimination forest
of $H^\Sigma_M(G)$ can be changed so that $u$ is a leaf child of
the deepest vertex in $X$, without increasing total depth or changing
the ancestor order among vertices of $G$.
\end{lemma}

\begin{proof}
Let $F$ be the given elimination forest and $y_X$ the deepest vertex
of $X$ in $F$. The edge $uy_X$ makes $u$
and $y_X$ comparable. Delete $u$, attaching its children to its
parent, or making them roots if $u$ was a root. Reinsert $u$ as a
leaf child of $y_X$, and call the resulting forest $F_{\rm new}$.
Deletion preserves the ancestor order among
the remaining vertices. Every neighbor of $u$ belongs to $X$ and
is ancestral to $y_X$, so reinsertion preserves every edge as well.

Write $d=\operatorname{depth}_F(u)$, $e=\operatorname{depth}_F(y_X)$,
and let $D_F(u)$ be the old proper descendants of $u$. Deleting $u$
reduces each of their depths by one; reinserting it as a leaf affects
no other depth. Its new depth $d'$ satisfies $d'\le d+|D_F(u)|$:
if $u$ was below $y_X$, then $d'=e+1\le d$; otherwise $d'=e$,
and the old path from just below $u$ through $y_X$ contains
$e-d$ vertices of $D_F(u)$. The single accounting identity
\[
 \totd(F_{\rm new})-\totd(F)=d'-d-|D_F(u)|\le0
\]
therefore proves the claim.
\end{proof}

\begin{restatedlemma}{Lemma~\ref{lem:np-total-depth-added-vertices}}
Every elimination forest of $H^\Sigma_M(G)$ can be transformed without
increasing total depth so that, for every $X\in\{A,B,C\}$, all
vertices of $T_X$ are leaf children of the deepest vertex in $X$.
Consequently,
\[
\totd(H^\Sigma_M(G))
 =\min_{F_G}\bigl(\totd(F_G)+MW(F_G)\bigr),\quad MW(G)
 \le\totd(H^\Sigma_M(G))
 \le MW(G)+\frac{N(N+1)}2.
\]
\end{restatedlemma}

\begin{proof}[Proof of Lemma~\ref{lem:np-total-depth-added-vertices}]
Apply Lemma~\ref{lem:added-vertex-move} to every added vertex.
Previously moved vertices remain leaf children of their respective
deepest vertices, so this gives the stated normalization.

Deleting the added leaves gives an elimination forest $F_G$ of
$\widehat G$, and their total contribution is $MW(F_G)$.
Conversely, every $F_G$ extends by these leaves to an elimination
forest of $H^\Sigma_M(G)$. This proves the identity.

The lower bound follows from $\totd(F_G)\geq0$ and
$W(F_G)\geq W(G)$. For the upper bound,
choose $F_G$ minimizing $W$. Its total depth equals $N$ plus the
number of strict ancestor pairs, so
$\totd(F_G)\leq N+\binom N2=N(N+1)/2$.
Substituting this forest in the minimization identity proves the
upper bound.
\end{proof}

Fix an elimination forest $F_G$ of $\widehat G$. For $v\in V(G)$, put
\[
\begin{gathered}
c(v)=|\{X\in\{A,B,C\}:v\text{ is ancestral to }y_X\}|,\\S_0=\{v\in V(G):c(v)=3\},\qquad
 S_1=\{v\in V(G):c(v)=2\}.
\end{gathered}
\]
These sets record the shared prefixes of the three clique paths.

\begin{lemma}[The same pair of descendants]
\label{lem:clique-path-pair}
If $S_1$ is nonempty, all its vertices are ancestral to the same two
vertices among $y_A,y_B,y_C$.
\end{lemma}

\begin{proof}
For two vertices of a rooted forest, their sets of descendants among
$y_A,y_B,y_C$ are disjoint or one contains the other. The sets for
vertices in $S_1$ have size two, so any two intersect and must be equal.
\end{proof}

When $S_1$ is nonempty, write $X$ for the part outside this common
pair and $Y,Z$ for the other parts. When $S_1$ is empty, choose any
part as $X$ and let $Y,Z$ be the other two parts.

For a choice of $X,Y,Z$ and sets $S_0,S_1$, write
$\Gamma_X(S_0,S_1)=2|S_0|+|S_1|$.
This is the cost of a feasible pair of covers. Equivalently,
$\operatorname{cost}_X(G)$ is the minimum of $\Gamma_X(S_0,S_1)$
over sets $S_0$ covering the $X$--$(Y\cup Z)$ edges and sets
$S_1\subseteq(Y\cup Z)\setminus S_0$ covering
$G[(Y\cup Z)\setminus S_0]$.

\begin{lemma}[Covers from the three paths]
\label{lem:clique-path-covers}
The set $S_0$ covers the $X$--$(Y\cup Z)$ edges of $G$, and
$S_1\subseteq Y\cup Z$ covers $G[(Y\cup Z)\setminus S_0]$.
Moreover,
\[
 W(F_G)=N+3+\Gamma_X(S_0,S_1)
 \ge N+3+\operatorname{cost}_X(G).
\]
\end{lemma}

\begin{proof}
Every vertex of a part is ancestral to its deepest vertex. A vertex
of $S_1$ has precisely $y_Y,y_Z$ as descendants among the three
deepest vertices, so it cannot belong to $X$. A vertex of
$X\setminus S_0$ is ancestral only to $y_X$, whereas a vertex of
$(Y\cup Z)\setminus S_0$ has a nonempty descendant set contained in
$\{y_Y,y_Z\}$. These sets are disjoint, so the vertices are
incomparable and cannot be adjacent. Hence $S_0$ covers the first
set of edges. Likewise, a $YZ$-edge outside $S_0$ with neither
endpoint in $S_1$ would join vertices ancestral only to $y_Y$ and
only to $y_Z$, which are incomparable. This proves the second cover
claim.

The sum of the three depths counts each vertex $v$ exactly $c(v)$
times. Every vertex belongs to its own part's path, so $c(v)\geq1$.
The definitions of $S_0,S_1$ therefore give
\[
 \sum_{X\in\{A,B,C\}}\operatorname{depth}_{F_G}(y_X)
 =\sum_{v\in V(G)}c(v)=N+2|S_0|+|S_1|.
\]
Adding the three terminal offsets gives the identity; the cover
claims give the inequality.
\end{proof}

\begin{restatedlemma}{Lemma~\ref{lem:np-total-depth-path-cost}}
For every tripartite graph $G$ with nonempty specified parts,
\[
 W(G)=N+3+\min\{\operatorname{cost}_A(G),
                 \operatorname{cost}_B(G),
                 \operatorname{cost}_C(G)\}.
\]
\end{restatedlemma}

\begin{proof}[Proof of Lemma~\ref{lem:np-total-depth-path-cost}]
The lower bound is
\[
 W(G)=\min_{F_G}W(F_G)
 \overset{\text{Lem.~\ref{lem:clique-path-covers}}}{\ge}
 N+3+\min_{X\in\{A,B,C\}}\operatorname{cost}_X(G).
\]

For the converse, choose a part $X$ attaining the minimum, and let
$Y,Z$ be the other two parts. Choose a cover $S_0$ attaining
$\operatorname{cost}_X(G)$ and a minimum-cardinality cover $S_1$
of $G[(Y\cup Z)\setminus S_0]$. Put $S_0$ on a common chain.
Below it, put the remaining $X$ vertices on one chain. On the other
branch put $S_1$ on a chain, followed by separate chains for the
remaining $Y$ and $Z$ vertices. Empty chains are omitted. Each part
lies on one ancestor path. Every cross-part edge either meets
$S_0$ or is a remaining $YZ$-edge meeting $S_1$, so this is an
elimination forest of $\widehat G$.

The path through $S_0$ to the $X$-chain and the two paths through
$S_0,S_1$ to the $Y$- and $Z$-chains contain every vertex of their
respective parts. If a final chain is empty, end its path at the
preceding shared chain. The depth of each part's deepest vertex is
at most its path length. Counting these three path lengths counts
every vertex once, every vertex of $S_0$ twice more, and every
vertex of $S_1$ once more. Hence
\[
 W(F_G)\le N+3+\Gamma_X(S_0,S_1).
\]
The last term equals $\operatorname{cost}_X(G)$, proving the upper bound.
\end{proof}

\begin{lemma}[Normalizing successive covers]
\label{lem:total-depth-cover-normalization}
Let $X,Y,Z$ be the specified parts of $G$. Suppose that $S_0$ covers
the $X$--$(Y\cup Z)$ edges and that
$S_1\subseteq(Y\cup Z)\setminus S_0$ covers
$G[(Y\cup Z)\setminus S_0]$. An independent block of vertices in
one part with identical neighborhoods can be made wholly contained
in $S_0$, wholly contained in $S_1$, or disjoint from both, without
increasing $\Gamma_X(S_0,S_1)$. The $S_1$ state is unavailable
for a block in $X$.

Moreover, if $v\in S_0$ and every neighbor of $v$ across the
$X$--$(Y\cup Z)$ cut already belongs to $S_0$, then $v$ can be
removed from $S_0$ without increasing the cost. If $v\in Y\cup Z$,
it suffices to add $v$ to $S_1$, decreasing the cost by one.
If $v\in X$, deleting it decreases the cost by two.
\end{lemma}

\begin{proof}
Give each vertex state $0$, $1$, or $2$ according as it belongs to
$S_0$, to $S_1$, or to neither; these states cost $2$, $1$, and $0$.
For a nonempty block $D$, choose a vertex of minimum state cost in
$D$, and give every vertex of $D$ its state. This cannot increase
$\Gamma_X$. It also preserves both covers: every edge from $D$
to a vertex outside $D$ now has the same endpoint states as the
existing edge from the chosen vertex to that neighbor. The first
cover requires a state-$0$ endpoint on each first-cut edge; the second
requires a state-$0$ or state-$1$ endpoint on each $YZ$-edge.
These conditions held on the representative edge, and $D$ has no
internal edges. A block in $X$ cannot acquire state $1$, because
none of its vertices had that state. An empty block needs no change.

For the last assertion, deleting $v$ preserves the first cover by
the stated condition. If $v\in X$, the residual graph is unchanged
and the cost decreases by two. Otherwise, every new residual edge
is incident with $v$, so moving $v$ to $S_1$ covers all such edges
and replaces cost two by cost one.
\end{proof}



\subsection{Sparse replacements on prescribed vertex sets}
\label{app:factor-expanders}

We use the Margulis--Gabber--Galil family from
Theorem~\ref{thm:auxiliary-expander-family} and the expander mixing
lemma (Lemma~\ref{lem:expander-mixing}).

Fix $\zeta>0$ and $\kappa\ge1$ as in Lemma~\ref{lem:factor-expanders}.
We map the vertices of a powered expander onto the prescribed sets
and verify the degree, no-hole, and connectedness guarantees.

Use $\lambda_0=5\sqrt2/8$ from
Theorem~\ref{thm:auxiliary-expander-family}. Fix an integer $h\ge1$
with $\lambda_0^h<\zeta/4$, and set
$d_{\zeta,\kappa}=\lceil4\kappa\rceil8^h$.
Use this same walk length for both constructions. For each integer
$\nu\ge1$, put $M=\nu^2$ and take
$\mathcal E=\mathcal G_\nu^{[h]}$ from that theorem.
It has degree $d=8^h$ and normalized operator norm at most
$\lambda_0^h$ on the zero-sum subspace, so its nonconstant
adjacency eigenvalue bound is $d\lambda_0^h$.
For $S,T\subseteq V(\mathcal E)$, Lemma~\ref{lem:expander-mixing} gives
\begin{equation}\label{eq:factor-expander-mixing}
 e_{\mathcal E}(S,T)
 \overset{\text{Lem.~\ref{lem:expander-mixing}}}{\geq}
 d\sqrt{|S||T|}\left(\frac{\sqrt{|S||T|}}{M}-\lambda_0^h\right).
\end{equation}
For an integer $t\ge1$, if $|S|,|T|\ge t>\lambda_0^hM$, then
$\sqrt{|S||T|}/M\ge t/M>\lambda_0^h$, so the right-hand side is
positive. Hence a step joins $S$ to $T$.

We accommodate arbitrary sizes by grouping auxiliary vertices and
assigning each group to a prescribed vertex. Auxiliary steps give
edges between distinct images; group sizes bound the output degrees.

For a map $\pi:V(\mathcal E)\to U$ and a subset $X\subseteq U$,
the \emph{preimage} is
$\pi^{-1}(X)=\{x\in V(\mathcal E):\pi(x)\in X\}$.
The map is \emph{onto} if every vertex of $U$ has at least one
preimage. When $1\le n=|U|\le M$, partition the $M$ auxiliary vertices
into $n$ groups. Write $M=n\lfloor M/n\rfloor+r$, where
$r$ is an integer with $0\leq r<n$. Give $r$ groups size $\lfloor M/n\rfloor+1$ and the
other groups size $\lfloor M/n\rfloor$. Since $M\geq n$, every group
is nonempty. Assign one group to each vertex of $U$. This gives
\begin{equation}\label{eq:factor-expander-preimages}
 \lfloor M/n\rfloor\le|\pi^{-1}(\{u\})|\le\lceil M/n\rceil
 \qquad(u\in U).
\end{equation}

Use the following two constructions.
\begin{enumerate}
 \item\label{constr:one-set-replacement} \emph{One prescribed set.}
 Let $U$ have size $n\ge1$, put $M=\lceil\sqrt n\rceil^2$, and
 take the $M$-vertex graph $\mathcal E$ above. Choose an onto map
 $\pi:V(\mathcal E)\to U$ with the group sizes just specified.
 Define the replacement $R_U$ by
 \[
 V(R_U)=U,
 \qquad E(R_U)=\{uv\in\tbinom{U}{2}:
 e_{\mathcal E}(\pi^{-1}(\{u\}),\pi^{-1}(\{v\}))>0\}.
 \]
 Thus equal images give no edge, and repeated pairs give a single edge.
 \item\label{constr:two-side-replacement} \emph{Two prescribed sides.}
 Let $A,B$ be disjoint sets of sizes $n_A,n_B\ge1$, put
 $M=\lceil\sqrt{\max\{n_A,n_B\}}\rceil^2$, and take the
 $M$-vertex graph $\mathcal E$ above. Choose onto maps
 $\pi_A:V(\mathcal E)\to A$ and
 $\pi_B:V(\mathcal E)\to B$, with preimage sizes
 $\lfloor M/n_A\rfloor$ or $\lceil M/n_A\rceil$ for the first
 map, and $\lfloor M/n_B\rfloor$ or $\lceil M/n_B\rceil$ for
 the second. Define the replacement $R_{A,B}$ by
 \[
 V(R_{A,B})=A\mathbin{\dot\cup}B,
 \qquad E(R_{A,B})=\{uv:u\in A,\ v\in B,
 e_{\mathcal E}(\pi_A^{-1}(\{u\}),\pi_B^{-1}(\{v\}))>0\}.
 \]
 This gives one edge for each distinct pair and no edge within
 either side. The two maps have the same
 domain, and their preimage groups may differ.
\end{enumerate}

A walk is \emph{closed} if its first and last vertices coincide.

\begin{restatedlemma}{Lemma~\ref{lem:factor-expanders}}
For every fixed $\zeta>0$ and $\kappa\geq1$, there is a constant
$d_{\zeta,\kappa}$ such that the following connected simple graphs,
each of maximum degree at most $d_{\zeta,\kappa}$, can be constructed
deterministically in polynomial time.
\begin{enumerate}
 \item For any specified nonempty finite set $U$, there is a graph on
 $U$ in which every two disjoint sets $X,Y\subseteq U$ with
 $|X|,|Y|\geq\lceil\zeta|U|\rceil$ have an edge between them.
 \item For any specified disjoint nonempty finite sets $A,B$ with
 $\max\{|A|,|B|\}/\min\{|A|,|B|\}\leq\kappa$, there is a bipartite
 graph on $A\mathbin{\dot\cup}B$ with parts $A,B$ in which every
 $X\subseteq A$, $Y\subseteq B$ with
 $|X|,|Y|\geq\lceil\zeta(|A|+|B|)\rceil$ have an edge between them.
\end{enumerate}
\end{restatedlemma}

\begin{proof}[Proof of Lemma~\ref{lem:factor-expanders}]
Each auxiliary vertex has $d=8^h$ outgoing and incoming steps,
because its adjacency matrix is symmetric with row sum $d$.
For any output vertex $v$, let $\rho$ be its map ($\pi$, $\pi_A$,
or $\pi_B$), and let $n_\rho$ be the size of that map's prescribed
set. Count outgoing steps from its preimage, except on side $B$,
where we count incoming steps. Each output neighbor uses at least
one such step, so the same bound covers all three cases:
\[
 \deg(v)\le8^h|\rho^{-1}(\{v\})|
 \overset{\eqref{eq:factor-expander-preimages}}{\leq}
 8^h\lceil M/n_\rho\rceil.
\]

In the one-set construction,
\begin{equation}\label{eq:factor-expander-order}
 M=\lceil\sqrt n\rceil^2\le(\sqrt n+1)^2\le4n.
\end{equation}
Hence $\lceil M/n\rceil\le4$. For the bipartite construction, the
same estimate with $n=\max\{n_A,n_B\}$ gives
$M\overset{\eqref{eq:factor-expander-order}}{\leq}4\max\{n_A,n_B\}$.
The side-size ratio
is at most $\kappa$, so
$M/n_A,M/n_B\le4\kappa$. Both constructions therefore have degree
at most $\lceil4\kappa\rceil8^h=d_{\zeta,\kappa}$.

We next verify the no-hole conditions at the thresholds in the
lemma. In the one-set case, put $t=\lceil\zeta n\rceil$. Then
$\lambda_0^hM<(\zeta/4)M
\overset{\eqref{eq:factor-expander-order}}{\leq}\zeta n\le t$.
In the bipartite case, put $t=\lceil\zeta(n_A+n_B)\rceil$. Then
$\lambda_0^hM<(\zeta/4)M
\overset{\eqref{eq:factor-expander-order}}{\leq}\zeta\max\{n_A,n_B\}\le t$.
We apply the mixing bound to preimages of $t$-sets. If no such admissible sets exist, the corresponding
no-hole condition is vacuous. Otherwise, any larger admissible sets
contain admissible $t$-element subsets.
For the one-set construction, take disjoint $t$-sets $X,Y\subseteq U$
and put $S=\pi^{-1}(X)$, $T=\pi^{-1}(Y)$.
For the bipartite construction, take $t$-sets $X\subseteq A$,
$Y\subseteq B$ and put $S=\pi_A^{-1}(X)$, $T=\pi_B^{-1}(Y)$.
In both cases surjectivity gives $|S|,|T|\ge t>\lambda_0^hM$,
so \eqref{eq:factor-expander-mixing} supplies a step from $S$ to $T$.
Its endpoint images give an edge between $X$ and $Y$. In the one-set
case the preimages are disjoint and their images are distinct.
In the bipartite case the preimages may overlap, and even a loop
step gives a valid edge because its images lie in the disjoint
sets $A,B$.

It remains to prove connectedness. If $\mathcal E$ had a nonempty
proper component $C$, then
$x=\mathbf1_C-(|C|/M)\mathbf1$ would be a nonzero zero-sum vector
with $A_{\mathcal E}x=dx$, contradicting its nonconstant eigenvalue
bound $d\lambda_0^h<d$. Thus $\mathcal E$ is connected.
It also has a loop at $(0,0)$: repeat $h$ times a loop of the
Margulis--Gabber--Galil graph at that vertex. This explicitly
supplies the odd closed walk needed below.

In the one-set construction, applying $\pi$ to an auxiliary walk
gives a sequence of equal or adjacent output vertices. Deleting
consecutive repetitions gives an output walk. Connectedness of
$\mathcal E$ and surjectivity of $\pi$ thus prove connectedness of
the output.

For the bipartite construction we need auxiliary walks of both
parities. From any $x\in V(\mathcal E)$, walk to $(0,0)$, traverse its loop, and
return along the reversed walk. Prepending this odd closed walk to a walk from
$x$ to $y\in V(\mathcal E)$ changes its parity. Map a walk
$x=x_0,x_1,\ldots,x_\ell=y$ to the alternating sequence
$\pi_A(x_0),\pi_B(x_1),\pi_A(x_2),\ldots$ in the output.
Each consecutive pair is an edge by construction, using the
reverse auxiliary step when the sequence goes from $B$ to $A$.
An even-length walk reaches $\pi_A(y)$, and an odd-length walk
reaches $\pi_B(y)$. Thus $\pi_A(x)$ is connected to both images of
every $y$, and surjectivity gives connectedness.

Both outputs are simple by their edge definitions.
The auxiliary order $M$ is at most four times the largest prescribed
set size. Theorem~\ref{thm:auxiliary-expander-family} constructs its
walk-power graph in polynomial time for fixed $h$. Map the endpoints
as specified above and sort the resulting pairs to merge repeated
edges, discarding equal images in the one-set construction. Since
$h$ depends only on $\zeta$, these operations take polynomial time
for the fixed parameters.
\end{proof}

